\documentclass[12pt,a4paper]{article}

% ======= Packages =======
\usepackage[margin=1in]{geometry}
\usepackage{setspace}
\usepackage{titlesec}
\usepackage{fancyhdr}
\usepackage{tocloft}
\usepackage{hyperref}
\usepackage{times}
\usepackage{tabularx}
% \usepackage[fontset=fandol]{ctex}
% \usepackage{kotex}
% ======= Section formatting =======
\titleformat{\section}{\bfseries\large}{\thesection}{1em}{}
\titleformat{\subsection}{\bfseries\normalsize}{\thesubsection}{1em}{}
\titleformat{\subsubsection}{\itshape\normalsize}{\thesubsubsection}{1em}{}

% ======= Header & Footer =======
\pagestyle{fancy}
\fancyhf{}
\fancyfoot[L]{Maximum Academic Press}
\fancyfoot[C]{Page \thepage\ of 14}
\fancyfoot[R]{Manuscript Template}
\renewcommand{\headrulewidth}{0pt}
\renewcommand{\footrulewidth}{0.4pt}

% ======= Document starts =======
\begin{document}

\begin{center}
    \textbf{\LARGE Preparing your paper}
\end{center}

\section*{Title Page}

The title page should include: a succinct title (less than 250 characters), a very concise running title (which should be less than 50 characters, for example, \textit{Plant Immune Mechanisms}).

The full names of all authors including their given names, the affiliations and location (including city, state, country and zip/post code) of all authors, the full name, and official email address(es) of the corresponding author(s).

\section*{Abstract}

A brief abstract up to 250 words should state the background and purpose of the research, methods (or rationale), main results and findings, and brief conclusions of the study. The abstract should not contain abbreviations or references and should not be structured.

% ======= Main Text Section =======
\newpage
\section*{Main Text}

The main text should contain the following twelve sections. If level-1, level-2, or level-3 headings are used, please ensure that the hierarchical numbering is clearly indicated. (Note: Do not add “," at the end of heading numbers).

\vspace{1em}
\tableofcontents
\newpage

% ======= Main Text Body =======
\section{Introduction}
Introduction should provide a background on the research topic, and a focused literature review that includes known and controversial findings, challenging issues, and the hypothesis on this research. At the end, briefly summarize the materials and methods as well as the research rationale. Note that the Introduction should not contain results or conclusions. 
\section{Materials and methods}
Materials and methods should be described in sufficient detail to allow the research work to be reliably reproduced in another laboratory, and to leave the readers in no doubt as to how the results were derived. Please also remember to include a section of “Statistical Analysis” where the statistical methods with enough details including threshold cut-off are provided to enable an appraisal of the work and an analysis of the data to reproduce the results. When the total word number exceeds the limit, it is suggested to check the redundance in Introduction, Results, and Discussion for removal rather than to briefly describe the methods without providing sufficient details. In a circumstance where it is difficult to compress other sections, some detailed methods, especially those that are routinely used in most laboratories, can be placed into the Electronic Supplemental Information.
\section{Results}

Results should be presented in a logical sequence in the text, tables and figures, usually matching the order as described in Materials and Methods. Repeated presentation of the same data in different forms should be avoided. The Results should not include a lengthy discussion.

\subsection*{For Tables and Figures:}

Authors should submit tables and figures with clear contents. Tables and Figures should be numbered consecutively in Arabic numerals as Table~1, 2, 3 and Figure~1, 2, 3.

\begin{itemize}
    \item Figures/tables should be supplied as single files for each.
    \item Letterings on figures should be in Times New Roman (non-bold), and the font size should be 8--10\,pt. Any letterings/data in bold should be explained in figure caption.
    \item The minimum line width for figures should be 0.4\,pt.
    \item The composite Figures should be labeled as Figure~1a, 1b, 1c, etc.\ with a, b, and c clearly labeled in each panel.
    \item Any color images should be saved in RGB color mode at 300\,dpi or higher resolution.
    \item Any mono line art should be saved in gray mode at 600\,dpi.
    \item TIFF, EPS, JPG and PDF files are the preferred format. It is recommended that you generate your figures in JPEG format before converting them to PDFs or uploading individual files. This will reduce the file sizes and enable you to obtain the figures close to the requirements.
    \item Tables should be in editable format. Please avoid using vertical rules and shading in table cells.
    \item Table and figure legends should be placed immediately above or below each, respectively. All designations in the tables or figures should use lower case alphabet, in the order of a, b, c.
\end{itemize}

\begin{enumerate}
    \item[\textbf{a.}] \textbf{For first submissions} (\textit{i.e.} not revised manuscripts), please incorporate the manuscript text, tables and figures into a single file (Microsoft Word or TeX/LaTeX) up to 20\,MB in size --- The figures, figure legends, tables, table captions should be inserted within the text at the appropriate positions--after the first place which mentioned it.

    \item[\textbf{b.}] \textbf{For the revised manuscript}, we suggest you submit the manuscript text, tables and figures, supplementary files separately.
\end{enumerate}

\section{Discussion}
The results should not be restated in the Discussion, but can be recapitulated to support or rebut existing hypotheses, conceived assumptions, or true facts stated in the introduction or literature. The Discussion should compare and relate the new or major findings in the Results to the existing body of knowledge in the field, in terms of improvement or further advance of current knowledge and technologies (Methods paper), and overall significance and contribution to the field as well as the future research focuses that stemmed from this study. 

\section{Conclusions}
A short conclusion of the study may be presented in a short Conclusions section, or at the end of the Discussion section.
\section{Electronic supplementary information}

\begin{enumerate}
    \item[\textbf{a.}] All lesser significant figures, tables, and procedures that support the main body of key results and conclusions in the text should be included as \textbf{Electronic Supplementary Information} and uploaded as a separate file (PDF) at the same time of manuscript submission.

    \item[\textbf{b.}] The availability of electronic supplementary information should be mentioned in a separate paragraph in the manuscript, placed immediately before the References, as: “\textbf{Electronic Supplementary Information}”.
\end{enumerate}

\section{Author contributions}

The roles and contributions of each author must be described in the subsequent manner: The authors confirm contribution to the paper as follows: study conception and design: Author X (last name and initials), Author Y (last name and initials); data collection: Author Y (last name and initials); analysis and interpretation of results: Author X (last name and initials), Author Y (last name and initials), Author Z (last name and initials); draft manuscript preparation: Author Y (last name and initials). Author Z (last name and initials). All authors reviewed the results and approved the final version of the manuscript. An author name can appear multiple times, and each author name must appear at least once.

For single authors, please use the following wording: The author confirms sole responsibility for the following: study conception and design, data collection, analysis and interpretation of results, and manuscript preparation.

\section{Acknowledgments}

Acknowledgments should include the funding agency and grant number which provided other resources. Individuals who have contributed to make the research possible, but not sufficiently qualified to be authors should also be included in this section.


\section{Ethical Statement (Optional.)}

Ethical statements ensure that all aspects of the research and publication process are conducted with honesty and transparency. The editorial office ensures that all ethical statements are appropriate and relevant to the study being reported. Manuscripts submitted without a clear and adequate ethical statement will be returned to the authors for revision and will not be considered for further review until the necessary ethical declaration is provided. If the manuscript does not involve this issue, state ``Not applicable.'' in this section.

Authors should include the following points in the ethical statement (if applicable) when submitting a paper:

\begin{itemize}
    \item \textbf{For investigations involving animal experimentation}
\end{itemize}

\textbf{(1) Ethics Review Approval:} The study must undergo review and approval by an ethics committee (such as the Institutional Animal Care and Use Committee, IACUC). The statement should include the project identification number, approval date, and the name of the ethics committee.

When collaborative research is performed internationally, the research protocol must be approved by research ethics committees in both the sponsoring and host countries.

\textbf{Template:} All procedures were reviewed and preapproved by the [ethical review committee or equivalent that has approved the use of animals], identification number: ***, approval date: ***.

\textbf{(2) Adherence to the 3Rs Principle:} The statement should clarify whether the research follows the ``Replacement, Reduction, and Refinement'' principles to minimize harm to animals.

\textbf{Template:} The research followed the ``Replacement, Reduction, and Refinement'' principles to minimize harm to animals.

\textbf{(3) Animal Welfare:} Provide details on the housing conditions, care, and pain management for the animals, ensuring that the impact on the animals is minimized during the experiment.

\textbf{Template:} This article provides details on the housing conditions, care, and pain management for the animals, ensuring that the impact on the animals is minimized during the experiment.

\textbf{(4) Informed Consent in Clinical Studies:} If the research involves pets or other domestic animals, it is necessary to state whether informed consent was obtained from the owners, and they should be made aware of any potential risks.

\textbf{Template:} Informed consent was obtained from the owners, with a clear explanation of any risks and the intention to publish the results. For more information, please see the author’s informed consent: [documentation].

\begin{itemize}
    \item \textbf{For investigations involving human participants}
\end{itemize}

\textbf{(1) Ethical Approval:} Research involving human participants, human materials, or human data must adhere to the Declaration of Helsinki and receive approval from a relevant ethics committee. All manuscripts reporting such research must include a statement with the name of the ethics committee and, where applicable, the reference number. If a study has received an exemption from ethics approval, this should also be noted in the manuscript, including the name of the committee that granted the exemption. Additional documentation supporting this may be requested by the Editor. Manuscripts may be rejected if deemed not conducted within an appropriate ethical framework. In exceptional cases, the Editor may reach out to the ethics committee for further clarification.

If an ethics committee has not approved a study before it starts, it’s typically not feasible to get approval afterward, and the paper might not be eligible for peer review. The editor has the authority to decide if the manuscript should go through peer review in these situations.

\textbf{Template:} The study was conducted in accordance with the Declaration of Helsinki, and the protocol was approved by the Ethics Committee of Gastrointestinal Tumors (Project identification code) on [date of approval].

\textbf{(2) New clinical tools and procedures:} When authors describe the implementation of a novel procedure or instrument in a clinical context, such as in a technical note or case report, they must provide a compelling justification in their manuscript. This justification should explain why the new method or device was considered superior to standard clinical practices for fulfilling the patient's clinical requirements. This requirement for justification is waived if the new procedure has already been authorized for clinical use at the authors’ facility. For any experimental application of a new technique or device where a distinct clinical benefit was not evident prior to the intervention, authors are anticipated to have secured approval from the ethics committee and obtained informed consent from the patients.

\textbf{(3) Consent to participate:} In all studies with human subjects, it is necessary to acquire informed consent from participants, or from their parents or legal guardians if the participants are under 16 years of age. This information must be clearly stated in the research manuscript. For papers that deal with vulnerable populations—such as unconscious patients, potentially incarcerated groups like prisoners, or situations where consent might not have been fully informed—the editors may exercise their discretion and possibly submit the manuscript for additional review by an editorial oversight committee. Consent must be secured for any data that could reveal personal identity, including biomedical, clinical, and biometric information.

If a donor has passed away before the study began and thus cannot provide consent but remains identifiable (such as with HeLa cells), the manuscript should acknowledge their origin and the absence of consent.

\textbf{Template:} ``Informed consent for participation was obtained from all subjects involved in the study.'' \textbf{OR} ``Informed consent for participation is not required as per local legislation [provide local legislation].'' \textbf{OR} ``Verbal informed consent was obtained from the participants. Verbal consent was obtained rather than written because [state the reason].'' \textbf{OR} ``Informed consent for publication was obtained from all identifiable human participants.'' \textbf{OR} ``This study utilized HeLa cells, which originated from a deceased donor. Since the donor passed away before the commencement of the research, it was not possible to obtain consent. We acknowledge the origin of these cells and have taken measures to ensure that the conduct of the research adheres to ethical standards and privacy protection.''

\begin{itemize}
    \item \textbf{Clinical Trial Registry}
\end{itemize}

Articles relying on clinical trials should quote the trial registration number at the end of the abstract. The journal also encourages the registration of such studies in a public trials registry prior to participants being enrolled.

\textbf{Template:} The study’s clinical trial registration number is Gastrointestinal Tumors registered with [the name and URL of the registry where your trial is registered]. Participant registration took place from [the participation registration period].

\begin{itemize}
    \item \textbf{Data Analysis and Public Datasets Research}
\end{itemize}

\textbf{Data Source Declaration:} If the research is based on public or secondary-use datasets, the source of the data, the method of acquisition, and compliance with ethical requirements should be clearly stated.

\textbf{Template:} The data used in this study were provided by [individual, group, or public source] and collected by/from [previously published sources or otherwise]. Therefore, no ethics committee approval was required for this study.

\subsection*{Example}

``All experimental procedures were conducted according to the guidelines provided by the Animal Care Committee. The Ethical Committee approved all experiments for the Experimental Use of Animals at Shenyang Agricultural University (LACUC Issue No: 2022030315, Approved Date: 2022.03.03).''~\cite{zhang2023}

\begin{thebibliography}{9}
\bibitem{zhang2023} Zhang J, Li D, Tian Q, Ding Y, Jiang H, \textit{et al.} (2023). The effect of kiwi berry (\textit{Actinidia arguta}) on preventing and alleviating loperamide-induced constipation. \textit{Food Innovation and Advances}, \textbf{2}(1):1–8. doi:10.48130/FIA-2023-0001
\end{thebibliography}

\section{Conflict of interest}

A statement must be included for all contributing authors who are involved in various kinds of conflicts. Examples of potential conflicts of interest include employment, consultancies, stock ownership, honoraria, paid expert testimony, patent applications/registrations, and grants or other resources which may potentially influence the decision derived from this research. If no conflict of interest is declared, the following statement should be declared in the manuscript: ``The authors declare that they have no conflict of interest.''

\section{Data availability}

A statement on data availability is required for all original articles which informs readers about the accessibility of research data linked to a paper and outlines the terms under which the data can be obtained. The statement should contain details about the location of the data that underpins the outcomes presented in the article. This may involve links to publicly archived datasets analyzed or produced during the research, where relevant. The term ``data'' refers to the fundamental dataset required to comprehend, reproduce, and expand upon the conclusions put forth in the article. We acknowledge that there are circumstances where making research data openly accessible might not be feasible, especially when it could jeopardize individual privacy. In such cases, the manuscript should still include a data availability statement, along with any prerequisites for accessing the data.

The table below contains template statements that you can use or adapt. A combination of more than one if required for multiple datasets.

\vspace{1em}

\begin{center}
\renewcommand{\arraystretch}{1.5}
\setlength{\tabcolsep}{6pt}
\begin{tabularx}{\textwidth}{|p{0.45\textwidth}|X|}
\hline
\textbf{Availability of data} & \textbf{Template for data availability statement} \\
\hline
Data derived from public domain resources. & The data that support the findings of this study are available in the [NAME] repository. These data were derived from the following resources available in the public domain: [list resources and URLs]. \\
\hline
The data are not publicly available due to [restrictions]. & The datasets generated during and/or analyzed during the current study are not publicly available due to [REASON WHY DATA ARE NOT PUBLIC] but are available from the corresponding author on reasonable request. \\
\hline
The data that support the findings of this study are available on request from the corresponding author, [initials]. & The datasets generated during and/or analyzed during the current study are available from the corresponding author on reasonable request. \\
\hline
Data sharing not applicable (no new data generated). & Data sharing is not applicable to this article as no datasets were generated or analyzed during the current study. \\
\hline
Data available within the article or its supplementary materials. & All data generated or analyzed during this study are included in this published article (and its supplementary information files). \\
\hline
\end{tabularx}
\end{center}

\section{References}

\subsection*{a. Citations in Text}
Example:

1.Finally, forests play a leading role in the global cycling of energy, carbon, water and nutrients$^{[4-7]}$ .

2.For the past millennia, the livelihood of humans has largely depended on forest resources, but these resources are not inexhaustible$^{[2,3]}$.


\subsection*{b. References List}

Please ensure that every reference cited in the text is also present in the reference list. Authors should ensure the accuracy of references, and are encouraged to use EndNote or other reference manager in writing the MS. Only one citation should appear within each reference number. If more than one reference is cited in one place, please separate the numbers by commas with one space.

List numbered references in the Literature Cited with numerals and period, without parentheses.
Include the following information (in this order):
Name(s) of author(s), last name first, followed by initials without periods. Include both (or all) initials for each author whenever they were included in the original article or book. Do not leave space between initials. Do not use a comma between surnames and initials—use commas only to separate different authors’ names. If a given reference has six or more authors, list the first five, then type “et al.” in the bibliography. If a reference has five or fewer authors, list them all.
 
Year of publication of the article or book, followed by a period, with no parentheses. If the article has recently been accepted for publication and is actually in press, list it in the Literature Cited section. Provide journal title and expected year of publication, plus volume and pages when known.
 
Title of article or chapter (see above for the policy of individual Annual Reviews journals regarding whether to include titles of articles or chapters).
 
Title of journal (full spelling) or book (not abbreviated unless part of a periodical series).

For a book reference, name(s) of editor(s).

Volume number, then a colon and inclusive page numbers, e.g., 3–10, 71–77, 100–109, 331–335, 1002–1003, 1198–1202, 1536–1538; if there is no volume number, inclusive page numbers preceded by a comma and “pp.”. The issue number can be included in parentheses immediately following the volume, if necessary, e.g., 10(4):123–130.










For a book reference, name(s) of editor(s).

Volume number, then a colon and inclusive page numbers, e.g., 3--10, 71--77, 100--103, 331--335, 1002--1003, 1198--1202, 1536--1538; if there is no volume number, inclusive page numbers preceded by a comma and ``pp.''. The issue number can be included in parentheses immediately following the volume, if necessary, e.g., 10(4):123--130.

For a book reference, place of publication, name of publisher, and edition, if necessary.

\subsection*{Example:}

\textbf{Journals:}

\begin{enumerate}
\renewcommand{\labelenumi}{[\theenumi].}
    \item Sun Q, Csorba T, Skourti-Stathaki K, Proudfoot NJ, Dean C, et al. 2013. R-loop stabilization represses antisense transcription at the \textit{Arabidopsis FLC} locus. \textit{Science} 340:619--621.
    \item Wei M, Liu Q, Wang Z, Yang J, Li W, et al. 2020. PuHox52-mediated hierarchical multilayered gene regulatory network promotes adventitious root formation in \textit{Populus ussuriensis}. \textit{Nature Communications} 228:1369--1385.
    \item Campobenedetto C, Mannino G, Beekwilder J, Contartese V, Karlova R, et al. 2021. The application of a biostimulant based on tannins affects root architecture and improves tolerance to salinity in tomato plants. \textit{Scientific Reports} 11:354.
\end{enumerate}

\textbf{Books:}

\begin{enumerate}
\renewcommand{\labelenumi}{[\theenumi].}
\setcounter{enumi}{3}
    \item Johnson FY, Liu XW, Velasco R, Smith M. (Eds.) 2019. \textit{Gene networks.} Vol. 15. Atlanta: Maximum Academic Press. 292 pp. \url{http://doi.org/10.xxxxxx.xxxx}
\end{enumerate}

\textbf{Book Chapter:}

\begin{enumerate}
\renewcommand{\labelenumi}{[\theenumi].}
\setcounter{enumi}{4}
    \item Myles S, Liu D. 2004. The apple genome sheds light on the evolution of rosaceae. In \textit{Fruits Research}, ed. Smith A. Atlanta: Maximum Academic Press. pp. 66--78. \url{http://doi.org/10.xxx.xxxxxx}
    \item Anu A, Sahni S, Kumar P, Prasad BD, White R, et al. 2016. Secondary metabolites in horticultural crops. In \textit{Plant Secondary Metabolites}, eds. Siddiqui MW, Prasad K. Vol.1. New York: Apple Academic Press. pp. 221--234. \url{https://doi.org/10.1201/9781315366326-15}
\end{enumerate}

\textbf{Patents:}

\begin{enumerate}
\renewcommand{\labelenumi}{[\theenumi].}
\setcounter{enumi}{6}
    \item Denecker J, Hoeberichts F, Muhlenbock P, Van Breusegem F, Van Der Kelen K. 2013. \textit{U.S. Means and methods for the reduction of photorespiration in crops.} \textit{WO Patent No. 2014147249A1.}
\end{enumerate}

\textbf{Theses and Dissertations:}

\begin{enumerate}
\renewcommand{\labelenumi}{[\theenumi].}
\setcounter{enumi}{7}
    \item Daniell D. 2005. \textit{Alternative oxidase is involved in the pathogenicity, development, and oxygen stress response of Botrytis cinerea.} Thesis. University of Washington, U.S. pp. 55--78.
\end{enumerate}

\textbf{Conference Proceedings:}

\begin{enumerate}
\renewcommand{\labelenumi}{[\theenumi].}
\setcounter{enumi}{8}
    \item Mohan Jain S. 2013. Mutation-assisted breeding for improving ornamental plants. \textit{Proc. XXII International Eucarpia Symposium, Section Ornamentals, Breeding for Beauty, Thuringia, 2013.} Chicago: AAA press. pp. 85--98. \url{http://doi.org/10.xxxxxx.xxxxxx}
\end{enumerate}


\end{document}
